\section{Mathematical analysis: what the evidence determines}
\label{sec:mathematics}

A measurement can answer some questions while leaving others open. We begin with a finite example in which two checks give identical results for programs with different behavior. This leads to a bound on the mistakes that any decision rule using those checks must make.

\subsection{Measurement classes and the finite error bound}
\label{sec:measurement-classes}

Begin with the retained finite Boolean challenge. It contained nine fixtures,
of which eight were executable and one did not compile. Among the eight
executable fixtures, three matched the specified output on every one of the
four input pairs and five did not. Yet all eight passed both the axiom-policy
gate and the printed-type gate. A gate is simply a check returning pass or
fail. Imagine now giving another decision rule only those two gate outputs,
with no code, candidate name, or behavior results. It sees exactly the same
input---two passes---for every fixture. A deterministic rule, meaning one
that always gives the same answer to the same input, must give all eight the
same decision.

If the rule calls all eight task-conforming, it makes five mistakes. If it
calls all eight nonconforming, it makes three. The better of these two choices
still makes three mistakes. No more elaborate arithmetic on the same two bits
can tell it which three programs were correct. The behavior check succeeds
in distinguishing them because it supplies another observation.

We can state this argument for any finite set. Let $\mathcal X$ be the set of
objects being checked and $n$ its number of objects, with $n>0$. A particular
object is called $x$. Let $Z(x)$ be the retained measurement for that object,
which may be a pair of gate outputs rather than a number. Let $T(x)$ be its
reference answer: $1$ for conformity to a specified task reference and $0$
for nonconformity. The letter $T$ here denotes a binary task answer, not a
waiting time. In this example the reference is the exhaustive four-input
check; this notation does not assert that an equally usable reference exists
for every historical task.

Group together objects with the same measurement value $z$. Within such a
group, write $n_{z0}$ for the number with reference answer $0$, and $n_{z1}$
for the number with reference answer $1$. These are counts of objects. Let
$\delta$ return either $0$ or $1$ using only $Z(x)$. It must answer
for every object. Define its finite-set error fraction as
\[
\operatorname{Err}_{\mathcal X}(\delta)
=\frac{\text{number of objects for which }\delta(Z(x))\ne T(x)}{n}.
\]
This is a dimensionless fraction of mistakes in this set, not yet a
population probability.

\begin{proposition}[Unavoidable mistakes when measurements merge different answers]
\label{prop:measurement}
Under the preceding assumptions, every deterministic decision rule using
only $Z$ satisfies
\[
\operatorname{Err}_{\mathcal X}(\delta)
\ \ge\ \frac{1}{n}\sum_z\min(n_{z0},n_{z1}).
\]
The sum runs over the distinct measurement values. Some such rule attains equality.
Consequently, zero mistakes are possible exactly when every measurement group
contains only one reference answer.
\end{proposition}

\begin{proof}
Fix one measurement group. Since the decision rule receives the same input
throughout that group, it must choose one answer for the whole group. Choosing
$0$ misclassifies its $n_{z1}$ conforming objects; choosing $1$ misclassifies
its $n_{z0}$ nonconforming objects. Therefore the smallest possible number of
mistakes in this group is the smaller count. Add those minimum counts over
the groups and divide by the total number of objects. Choosing the majority
answer separately in every group achieves the bound. A group's minimum is
zero precisely when one of its two counts is zero.
\end{proof}

For the eight executable fixtures there is just one group, with $n_{z0}=5$
and $n_{z1}=3$. Substitution gives
\[
\text{smallest possible error fraction}=\frac{\min(5,3)}8=\frac38.
\]
The interpretation ``both gates passed, therefore the task was met''
instead has error fraction $5/8$. Thus $5/8$ describes that particular
accept-all interpretation, while $3/8$ is a limitation of \emph{every forced
deterministic binary rule restricted to the same two gate outputs}. Neither
is a historical error-rate estimate, and neither has the denominator needed
for defect-detection sensitivity. The noncompiling ninth fixture is not
silently added to this eight-object calculation.

Allowing the rule to abstain changes the question: it could avoid all mistakes
by refusing every decision, so we must also report the fraction it answers.
Randomization cannot lower the minimum \emph{expected} error here. If it
chooses $1$ with probability $r$, where $0\le r\le1$, its expected number of
mistakes in the single group is $5r+3(1-r)=3+2r$, whose minimum is $3$.
Individual randomized runs need not equal their expectation. These remarks
explain why the assumptions ``deterministic'' and ``answers every object''
appear explicitly in the proposition.

\subsection{Why acceptance records alone do not identify accuracy}
\label{subsec:accuracy}

The previous result assumed that reference answers were available and showed
what a coarse measurement loses. A different difficulty arises when reference
answers are absent. Suppose, for illustration, that the only retained outcomes
for four records are accept, reject, accept, reject. If all four reference
answers agreed with those outcomes, the recorded decisions would be fully
accurate. If every reference answer disagreed, they would be entirely
inaccurate. Both hypothetical lists are compatible with the four observed
decisions \emph{when no additional evidence constrains the reference answers}.
We have not changed mathematical truth; we have exhibited what the supplied
records alone fail to settle.

In symbols, let $A$ denote a recorded binary acceptance decision and $T$ a
reference answer. Observing $A$ without information that constrains $T$ cannot
choose between the possibilities $T=A$ and $T=1-A$. Their fractions of correct
decisions differ. This is called \emph{nonidentification}: different answers
to the research question remain consistent with what was observed. In a
four-record example the possible accuracy fractions are multiples of $1/4$,
not every real number between zero and one. Source analysis, an explicit
task reference, or independent adjudication may eliminate these possibilities;
the limitation concerns the stated observation scheme.

Historical labels also need a clear coding rule. Coding only ``pass'' as $1$
and all other labels as $0$ measures recorded acceptance versus nonacceptance.
If the other labels include partial or inconclusive decisions, this coding
does not automatically measure reviewer error. An accuracy calculation among
decisive pass/fail reviews would use a different denominator and report the
omitted fraction separately.

\subsection{Stopping, observation eligibility, and path exit}
\label{sec:process-theory}

\subsubsection{From a reconstructed segment to an observed path}
A full segment connects records through the high-confidence candidate-revision edges of Section~\ref{subsec:segment-rule}. The observed path is the portion followed under the stopping rule of Section~\ref{subsec:failure-cohort}. One step is a transition between submitted versions, rather than an internal edit, a model call, or a fixed amount of time. The path is reconstructed from saved records, and one task can have several separate paths.

Suppose the recorded sequence is fail, fail, pass. Starting at the first failure, there are two steps, and observation stops at the pass. Three candidate records have thus supplied two transitions. The chain might instead end while its latest review is still nonpassing; we call this a \emph{path exit}. The task could continue elsewhere, so exit does not imply permanent failure. Section~\ref{sec:process} applies this observation rule to the history. We first examine its consequences in a simple probability model.

Throughout the model, success means a recorded pass. A \emph{nonpass} includes every other label, including partial or inconclusive judgments, whereas a failure-origin path must begin with an explicit failure label. These conventions concern the review record, rather than independently certified task conformity.

\subsubsection{Constant conditional pass probability and predictable observation}
Consider repeated tosses of a coin with head probability $p$, where $0<p<1$, stopping at the first head. If heads represents a recorded pass, this gives a model of successive submissions with a constant pass probability. Real submissions may improve and their review conditions may change; those possibilities are set aside in this illustrative model.

To express exactly what the calculation needs, let $j=1,2,\ldots$ denote a possible step. Write $H_{j-1}$ for all information available \emph{before the result of step $j$}: earlier candidates, earlier review outcomes, and any observation decision already made. Define the observation and acceptance indicators, and their cumulative counts:
\begin{center}
\begin{tabularx}{\linewidth}{lY}
\toprule
Symbol & Meaning for a single path \\
\midrule
$B_j$ & 1 if no previous step has passed and step $j$ is to be observed; 0 otherwise. \\
$A_j$ & 1 if the review at that step records a pass; 0 if it does not. Set it to 0 when $B_j=0$, merely to complete the notation. \\
$N_k=\sum_{j=1}^k B_j$ & Number of eligible steps actually observed through step $k$. \\
$D_k=\sum_{j=1}^k B_jA_j$ & Number of recorded passes among those steps. With first-pass stopping this is at most 1. \\
\bottomrule
\end{tabularx}
\end{center}

Saying that $B_j$ is known from $H_{j-1}$ means that we decide whether this observation belongs in the analysis without first seeing its outcome. This property is often called \emph{predictable observation}. It is not a prediction that the next candidate will pass. For example, ``continue after a failure and stop after a pass'' uses information already available. ``Retain this row only if the next review looks promising'' can use information that was not yet available and need not satisfy this condition.

For an eligible step, the constant-probability assumption is
\begin{equation}
 P(A_j=1\mid H_{j-1})=p\qquad\text{whenever }B_j=1.
 \label{eq:constant-conditional-pass}
\end{equation}
This requires the same conditional pass probability after every eligible past history; equality of pooled pass proportions is weaker. Equivalently, $E(A_j\mid H_{j-1})=p$ on eligible histories.

\begin{proposition}[A constant conditional probability under an explicit observation rule]
\label{prop:stopping}
Suppose observation eligibility $B_j$ is known before the step-$j$ outcome, and suppose equation~\eqref{eq:constant-conditional-pass} holds for every eligible history. For each fixed finite number $k$,
\begin{equation}
 E(D_k-pN_k)=0,\qquad\text{and hence}\qquad E(D_k)=pE(N_k).
 \label{eq:stopped-count-identity}
\end{equation}
If, in addition, the first step is observed and every nonpassing step is followed by another observed step until the first pass, the first-pass step $T$ satisfies
\begin{equation}
 P(T=j)=(1-p)^{j-1}p,\qquad j=1,2,\ldots.
 \label{eq:geometric-pass}
\end{equation}
\end{proposition}

\begin{proof}
At an eligible step the expected value of $A_j-p$, given the past, is $p-p=0$. At an ineligible step $B_j(A_j-p)=0$. Since eligibility is already known, both cases give
\[
 E\{B_j(A_j-p)\mid H_{j-1}\}=0.
\]
Averaging over the possible past histories still gives zero. Add this identity for the fixed set of steps $1,\ldots,k$. The sum inside the expectation is precisely $D_k-pN_k$, proving equation~\eqref{eq:stopped-count-identity}.

With observation continuing after every nonpass, a first pass at step $j$ requires $j-1$ consecutive nonpasses followed by a pass. Multiplying the successive conditional probabilities gives
\[
 \underbrace{(1-p)\cdots(1-p)}_{j-1\text{ nonpasses}}p.
\]
This is equation~\eqref{eq:geometric-pass}. It tends to put more weight on early steps because a later first pass requires surviving all the earlier nonpasses. The calculation uses the stated conditional probabilities; a separate independence assumption is not needed in addition to that full-history condition.
\end{proof}

Equation~\eqref{eq:geometric-pass} gives a geometric first-pass step $T$, distinct from the task-conformity reference $T(x)$. For $p=0.2$, the probabilities of first passing at steps 1, 2, and 3 are $0.2$, $0.8(0.2)=0.16$, and $0.8^2(0.2)=0.128$. Even though fewer paths remain eligible at each step, the chance of passing \emph{among paths still eligible} remains $0.2$ in this example. Thus a shrinking set of eligible paths does not, by itself, refute a constant-probability model. It also does not establish that such a model fits the historical archive.

\subsubsection{Martingale formulation}
\label{subsec:martingale}
Some earlier analyses used the word \emph{martingale}. Here it describes the running discrepancy
\[
 M_k=D_k-pN_k=\sum_{j=1}^k B_j(A_j-p).
\]
Given the available past, its next change has mean zero: $E(M_k\mid H_{k-1})=M_{k-1}$. That is the relevant mathematical property, already proved above. A filtration is simply the increasing family of information sets represented here by the histories $H_0,H_1,\ldots$. The martingale property is therefore another way to state the zero-mean increment calculation.

\subsubsection{Averages with a random denominator}
Equation~\eqref{eq:stopped-count-identity} concerns expected \emph{counts}. It does not say that $E(D_k/N_k)=p$, because the number of observed steps can itself depend on earlier outcomes. A short coin example makes the distinction visible. Take $p=1/2$, stop at the first pass, and observe at most two steps:
\begin{center}
\begin{tabular}{lrrrr}
\toprule
Observed sequence & Probability & $D_2$ & $N_2$ & $D_2/N_2$ \\
\midrule
Pass & $1/2$ & 1 & 1 & 1 \\
Nonpass, pass & $1/4$ & 1 & 2 & $1/2$ \\
Nonpass, nonpass & $1/4$ & 0 & 2 & 0 \\
\bottomrule
\end{tabular}
\end{center}
Here $E(D_2)=3/4$ and $E(N_2)=3/2$, so $E(D_2)/E(N_2)=1/2$. But the average of the final path proportions is $E(D_2/N_2)=1/2+1/8=5/8$. Fast successes contribute a ratio of 1 after only one observation. Exchanging ``take a ratio'' and ``take an expectation'' therefore changes the answer.

If observation instead continues without a step limit, use $D$ and $N$ for the \emph{total} pass count and observed-step count at the eventual stopping point; these differ from $D_k,N_k$, which count only through a fixed finite step $k$. Since $P(T>k)=(1-p)^k$ tends to zero, a first pass occurs with probability 1. Consequently, $D=1$ and $N=T$ with probability 1, and that probability-zero exception does not affect the following expectation. The same distinction has the exact expression
\[
 E(D/N)=\sum_{j=1}^{\infty}\frac{p(1-p)^{j-1}}{j}
 =\frac{-p\log p}{1-p}>p.
\]
Integrating the geometric series $\sum_{r=0}^{\infty}x^r=1/(1-x)$ from 0 to $1-p$ gives $\sum_{j\ge1}(1-p)^j/j=-\log p$. Also $-\log p=\int_p^1(1/u)\,du>1-p$ for $0<p<1$. This is a property of the illustrative model, not a correction formula applied to the historical data.

\subsubsection{Exit and selection of successful paths}
The no-exit coin model assumes that every nonpass is followed by another observation. Now suppose that after each nonpass a second independent coin determines whether the path continues. Let $a$ be its continuation probability, with $0\le a\le1$. All these continuation coins and acceptance coins are independent in this illustrative construction. The events after an eligible step are then:
\begin{center}
\begin{tabularx}{\linewidth}{Yl}
\toprule
What happens & Probability \\
\midrule
Record a pass and stop & $p$ \\
Record a nonpass and exit the path & $(1-p)(1-a)$ \\
Record a nonpass and observe another step & $(1-p)a$ \\
\bottomrule
\end{tabularx}
\end{center}
These probabilities add to 1. Give the third one the short name $b=(1-p)a$: it is the chance of both not passing now \emph{and} reaching the next observation. It is not the probability of a pass.

For a pass to be observed at step $j$, the first $j-1$ steps must each produce the third event, and step $j$ must pass. Hence
\[
 P(\text{pass at step }j\text{ before exit})=b^{j-1}p.
\]
The events for different $j$ cannot both occur. Summing their probabilities and using $1+b+b^2+\cdots=1/(1-b)$ gives
\begin{equation}
 s:=P(\text{pass before exit})=\frac{p}{1-b}.
 \label{eq:pass-before-exit}
\end{equation}
Here $s$ is a \emph{whole-path} probability; $p$ is a probability for one eligible step. They have different denominators and answer different questions.

Finally, suppose we discard all paths that exit without a pass and report only successful paths. If $J$ is the step of their recorded pass, ordinary conditional probability gives
\begin{equation}
 P(J=j\mid\text{pass before exit})
 =\frac{p b^{j-1}}{s}=(1-b)b^{j-1}.
 \label{eq:selected-success-time}
\end{equation}
This is another geometric distribution, but its parameter is $1-b=p+(1-p)(1-a)$, which exceeds $p$ if $a<1$. Successful-path selection favors shorter histories: to be a long successful path, a record must survive additional opportunities to exit. The faster-looking distribution does not demonstrate a higher per-step pass probability.

For a concrete identification problem, consider two possible mechanisms:
\begin{center}
\begin{tabular}{lrrrr}
\toprule
Mechanism & $p$ & $a$ & $b=(1-p)a$ & $s=p/(1-b)$ \\
\midrule
I & 0.2 & 0.5 & 0.4 & $1/3$ \\
II & 0.5 & 0.8 & 0.4 & $5/6$ \\
\bottomrule
\end{tabular}
\end{center}
Both give exactly the same successful-path step probabilities: $0.6$, $0.24$, $0.096$, and so on. Yet their eligible-step pass probabilities are $0.2$ and $0.5$. Looking only at successful-path lengths cannot distinguish these mechanisms. The overall pass fraction would distinguish them: if both $s$ and $b$ were known, equation~\eqref{eq:pass-before-exit} would give $p=s(1-b)$. The ambiguity comes from what was retained, not from a failure of probability arithmetic.

The three quantities now have distinct meanings: the pass rate at an eligible step, the outcome of a fixed path, and the length of a path selected because it passed. Applying the model to historical records would require checking its observation assumptions. A retrospective edge classification can use the next candidate or review, so reproducing that classification does not establish the before-outcome condition in Proposition~\ref{prop:stopping}. The empirical analysis below is consequently restricted to the retained paths and their observed associations.

\subsection{Partial tables and logical constraints}
\label{sec:feasible-tables}

Return to the four-cell table in Section~\ref{sec:framework}. Sometimes we
have measured only the second criterion's column. We may nevertheless know
something about the first column. For example, if one axiom policy allows a
subset of the axioms allowed by the other, passing the more restrictive policy
logically implies passing the more permissive one for the same extracted
declaration. The direction of this implication must be verified; calling a
policy ``new'' does not establish that it is stricter.

Assume for the following \emph{illustrative cases} that both scores are binary
and that acceptance under $R_1$ implies acceptance under $R_0$. In symbols,
\[
q_{i1}\le q_{i0}\qquad\text{for each candidate }i\in\{0,1\}.
\]
For zero-one scores this inequality means simply that a right-column $1$
requires a left-column $1$. It does not impose the reverse implication.

If the measured right column is $(1,1)$, the left column must be $(1,1)$ too.
The only permitted full table consists of four ones, so all six differences
are already zero. Actually running the two left-column checks could verify
the assumption or reveal a contradiction, but cannot give these differences
another value while the assumption and observations remain valid.

If instead the measured right column is $(0,0)$, each left-column value may
still be either zero or one. The four permitted completions are listed below.
Write $d=q_{10}-q_{00}$ for the candidate difference under $R_0$. Since the
right-column difference is zero, the definitions give $I=-d$ and
$\phi_C=d/2$.

\begin{center}
\begin{tabular}{ccccrrr}
\toprule
$q_{00}$ & $q_{10}$ & $q_{01}$ & $q_{11}$ & $d$ & $I$ & $\phi_C$\\
\midrule
0 & 0 & 0 & 0 & 0 & 0 & 0\\
0 & 1 & 0 & 0 & 1 & $-1$ & $1/2$\\
1 & 0 & 0 & 0 & $-1$ & 1 & $-1/2$\\
1 & 1 & 0 & 0 & 0 & 0 & 0\\
\bottomrule
\end{tabular}
\end{center}

Thus $d$ can take exactly the values $-1,0,1$. Its lower and upper bounds are
$-1$ and $1$, but $0.3$ is not possible in this binary example. Likewise,
$\phi_C$ can be $-1/2,0,1/2$. The output choices are linked: we cannot take
$d=1$ from one row and $I=1$ from another and claim they describe one table.
They must come from the same completion.

\subsection{Identification of a linear comparison}
\label{subsec:linear-identification}

Call $\mathcal F$ the collection of all full tables that agree with the
observed cells, the score's permitted values, and the verified logical rules.
Represent a table by the vector $q=(q_{00},q_{10},q_{01},q_{11})^{\top}$. Here $\mathcal F$ names a collection of
possible tables, not a probability distribution over them. Assume for now
that at least one table is possible.

Let $L(q)$ be a comparison calculated from a full table. For the comparisons
used here it is a weighted sum
\[
L(q)=a_{00}q_{00}+a_{10}q_{10}+a_{01}q_{01}+a_{11}q_{11},
\]
where $a=(a_{00},a_{10},a_{01},a_{11})^{\top}$ is a fixed weight vector, so $L(q)=a^{\top}q$. For example, the
candidate difference under $R_0$ uses weights $(-1,1,0,0)$. The output remains
in score points. Saying $L$ is \emph{uniquely determined}, or
\emph{point-identified}, means that every permitted table gives the same
answer.

\begin{proposition}[When the available information determines a comparison]
\label{prop:completion}
The comparison $L$ is uniquely determined exactly when $L(q)=L(q')$ for
every two permitted tables $q$ and $q'$ in $\mathcal F$. For a weighted sum,
this condition can be checked by verifying
\[
a_{00}(q_{00}-q'_{00})+a_{10}(q_{10}-q'_{10})
+a_{01}(q_{01}-q'_{01})+a_{11}(q_{11}-q'_{11})=0
\]
for every such pair. Further valid evidence can only remove permitted tables.
If it leaves at least one table, the greatest lower bound of the answers
cannot decrease and their least upper bound cannot increase. For a finite
list these bounds are simply its smallest and largest values.
\end{proposition}

\begin{proof}
If two permitted tables give different answers, the observations and rules
allow both answers, so there is no unique answer. If every permitted table
gives the same answer, that answer is forced. For a weighted sum, subtracting
$L(q')$ from $L(q)$ gives the displayed expression. Finally, removing choices
from a list cannot produce a new value below its previous minimum or above
its previous maximum. The same statement uses greatest lower and least upper
bounds when the permitted scores form an infinite set.
\end{proof}

In the four-row example, observing $q_{00}=0$ removes the last two rows. The
possible values of $d$ shrink from $\{-1,0,1\}$ to $\{0,1\}$. Also observing
$q_{10}=1$ leaves only the second row and determines $d=1$. This is additional
information. By contrast, computing $\phi_C=d/2$ from each of the original
four rows removes none of them. A useful summary need not be a new
observation. The two-factor allocation can aid explanation without supplying
a missing cell or a missing task reference.

If a measurement contradicts a claimed rule, there may be no permitted table
at all. For example, an assertion that the two criteria always give the same
score for a candidate conflicts with measured scores $0$ and $1$ for that
candidate. This calls for an inconsistency report, not an invented completion.
The release's consistency guard rejects this specific conflict. Invalid or
inapplicable comparisons likewise need their meaning repaired before we
enumerate numerical possibilities.

For meaningful but unknown cells, the released numerical core allows values
between $0$ and $1$ and calculates bounds over separate per-cell intervals.
Each cell receives its own permitted range, and the calculation allows any
combination of values from those ranges without expressing constraints
between cells. This does not assume that random variables are independent
in the probabilistic sense. These are bounds on missing information, not statistical confidence
intervals. If the score is actually binary, the exact attainable values may
be discrete as above. If additional relations connect the cells, separate
per-cell intervals may also allow impossible combinations and yield wider bounds than
necessary. The baseline code makes specific equality and binary-implication
deductions; it does not implement a general solver for every collection
$\mathcal F$. The proposition explains the reasoning principle rather than
claiming an unimplemented capability.

This distinction also clarifies the historical comparison: verified policy
implications already determine the accepted cases' missing policy cells,
while other configurations can benefit from measuring those cells. Equal
availability of the six specified diagnostics with and without an allocation
does not establish that allocation is universally useless. It reflects what
information those particular questions require.
